\section{Talking Father’s Day Blues}

I had time this morning to watch some news shows. The same topics come up day after day, at least until they fade from memory. For example, does anyone even care about North Korea any more? With three cable channels broadcasting news 24×7, you would think there would be some deep understanding of the issues. Nevertheless, I gained no such understanding today. It's not that the newsmen fail to ask the politicians the hard questions; they don't even ask them the obvious questions.

\paragraph{Syria}


Today, Syria was the hot topic since the USA has committed to arm and support the ``rebels''. We hear that 90,000 Syrians have been killed, but surely Assad and his men are not responsible for all of them. It seems to me that if the rebels were prevented from acquiring weapons, the rebellion would simply peter out. In the meantime, the USA is supporting a UN treaty that would disarm the American population. Why not apply it to the Syrians?

It was never made clear why Assad must go. The Christian bishops have been supporting Assad and the rebels have connections to Al Qaeda, the avowed enemy of the USA. That should count for something. The rebellion has become an international proxy war, much like the Spanish Civil War of the 1930s.

Soon after Obama's famous ``Red Line'', it has been proved that Syria is using chemical weapons. Very convenient, just as Iraq has weapons of mass destruction. Strangely enough, while it is against international law to use poison gas in a war, it is not so when used against a nation's own population. So nobody is criticizing Turkey for using tear gas. Is the war in Syria subject to international law or is it a civil disturbance as in Turkey or even in France last month? I am not an expert in law, but the question should have been asked.

The accusation is often made that Assad is killing his own people. Every government will use force to maintain power and control. But that standard, how many of ``his own people'' did Abraham Lincoln kill? Look it up ... he is considered America's greatest president. More recently, how many of ``his own people'' did Bill Clinton kill in Waco Texas?

\paragraph{The New Canal}


China wants to build a canal through Nicaragua to handle the new generation of super large tankers that cannot pass through Panama. When completed, they will be able to ship their junk to Walmart more efficiently.

\paragraph{The NSA}


The massive surveillance on American citizens by the National Security Administration was another topic, actually quite divisive. Americans always made fun of the ``secret police'' among the Nazis or in the Soviet Union. However, there is insufficient detachment to notice the parallels.

A defender of the practice pointed out that the phone company keeps the same metadata records. But they keep them for a limited time for billing purposes and the convenience of the customers. More importantly, they don't mine the data to discover who is contacting prostitutes, having an affair, and so on.

\paragraph{Population Movement}


Population movement has been going on since the beginning of time. The Romans were replaced by Germanic tribes. More recently, this has occurred in the Middle East. Besides the obvious example, in the 1960s Britain deliberately changed the demographics of Bahrain in order to dilute the Iranian population. What goes around, comes around.

The immigration bill in the USA is a similar example of deliberate demographic change. A Republican senator claimed the bill had to be passed or else they would never again win a national election. He offered no proof or explanation. The hair of the dog\footnote{\url{http://en.wikipedia.org/wiki/Hair\_of\_the\_dog}} is not always an effective cure.

I think Europeans are simply bored as they have perhaps actualized all their possibilities. There has been no great saint in centuries. Art, poetry, etc. have likewise reached a dead end. The best among us have become nihilists with nothing to offer except suicide or ``riding the tiger''. A Dutch woman\footnote{\url{http://www.youtube.com/watch?v=wjk0aWTeOh4}} has to pretend to be Hispanic.

\paragraph{Creating an Order}


Someone mentioned a Vedantic order that will accept ``anyone''. It can't be much of an order. The problem is not the lack of men capable of leading an order, but rather the shortage of men to enter it. Men are too presumptuous today.

\paragraph{Freedom and the Spiritual Vitality of Man}


\begin{quotex}
All of man's freedoms serve to make the material and spiritual life worthy of his station accessible to him. It is for this reason that he is in need of freedom, and must free himself from the bonds of his lower emotions and concupiscence in order to obtain this freedom. Everyone is susceptible to and a captive of acting on their lower urges and desires, the only exceptions being those in the ranks of the righteous.

Your souls are in bondage to your deeds; so free your souls by repenting of your sins and asking forgiveness of God. Thus man's real freedom is achieved by releasing himself from the bonds of past sins through repentance and forgiveness. It is only this spiritual law and emotional freedom that can lay the requisite groundwork for authentic and free human existence and which is the freeing of one's neck from the burden of sin, which results in the true experience of freedom.

The life of the spirit is more important than that of the body, and inner spiritual freedom is more important than outer social freedom and is in fact its cause and origin. The soul is only truly alive when it has freed itself from bondage to its lust and anger.
\end{quotex}


Are you thinking I have quoted some Catholic saint or Doctor of the Church? No, the author is \textbf{Abdollah Javadi-e Amoli}, a Shi'a theologian. It's odd that an external enemy can be a spiritual friend, while some who claim to be close to us are in fact spiritual enemies. The latter are the worst kinds of enemy.

\url{https://www.youtube.com/watch?v=wjk0aWTeOh4}

\flrightit{Posted on 2013-06-16 by Cologero}

\begin{center}* * *\end{center}

\begin{footnotesize}
\begin{sffamily}
\texttt{August on 2013-06-17 at 09:32 said:}

I thought as much about this Vedantic order. Here are a few interesting anecdotes about it.

I once encountered a white adherent of this ``order'' in Australia, in orange robes and tonsured. He approached us, seeking change, in a pub. At the time I was a somewhat ignorant but thoroughly inquisitive and skeptical youth, and I asked him what his people do with the money, without handing over a coin. He seemed to summon and concentrate what I felt to be a malefic charge in his eyes, glaring at me with no response. After a few seconds my companion, a rather ordinary soul, placed some coins in his pan and he went away.

Some time later, I was a slightly older but still skeptical and obstinate youth, when a similar projection was directed at me by another white fellow, this time conventionally clothed. He invited me to purchase a sticker with some kind of smiling sun design on it ( bayimg.com/KaNMCaaEE ). I declined the offer, when he became annoyed, and challenged me, sloppily but aggressively suggesting that I was very lucky and should share my happiness, or some such insight, though never telling me anything about his efforts. Later I found out that these stickers are apparently connected to the same ``order'', and depict the Lord Krishna, of whom this man perhaps fancied himself a representative. I am glad that he did not decide to emulate Arjuna.

Perhaps both of these chaps were ultimately vexed by the reminder that they could not effect a rudimentary psychic takeover of every Westerner they engaged, let alone a Hinduesque spiritual renaissance.

\hfill

\texttt{Jason-Adam on 2013-06-17 at 13:09 said:}

Is the order you are discussing ISKCON ?

I am one of those tiger riders, but not by choice. To simply put it, being a traditionalist Catholic man in Gay York, what else can I do ?

Reading this post reminds me of Spengler's winter of the west. As a western man it saddens me and I wonder why I chose to be born in this time. I have hope though for the future of the West because precisely I as a strong Western man would not have chosen to be born in a time where I could be of no use to anyone, there will come a new era in my lifetime I am sure where I and others like ourselves will be needed and able to act. In the meantime Cologero is doing the world a service by being a voice of truth against the modernists, counterinitiates, and demons.

\hfill

\texttt{Constantine Aetos on 2013-06-18 at 00:27 said:}

Two short responses: 1) since when the Dutch were ever a nation 2) it's not that men are presumptuous, they've become to feminized for God.

\hfill

\texttt{John Morgan on 2013-06-19 at 07:06 said:}

The original reference that was made in a prior comment thread to an article about a Western Vedantic order was to the article ``Vedic Veersa'' by Patrick Boch, which was published in the second issue of Arktos' journal, The Initiate, and which was indeed about ISKCON. I myself was involved with ISKCON for three years, although I no longer am. To say that it ``accepts anyone'' is not accurate. As with most other religions, of course anyone is free to come to the temples, observe the services, learn the philosophy and do the basic practices. But there are two stages of initiation in Gaudiya Vaishnavism, harinamadiksa and brahmana, and one must study and practice for years under the guidance of a spiritual master before one can even become eligible to apply for the testing that is a prerequisite for first initiation (and not everyone passes the tests). The second stage is even more difficult. Whether one finds the philosophy and practices palatable for a Westerner is debatable (which is why I ultimately left), but it did satisfy me as possessing all the qualifications of a traditional, initiatic path, and its founder, A. C. Bhaktivedanta Swami Prabhupada, definitely qualifies as one of the greatest spiritual leaders of the twentieth century, if one considers his biography, traditional teachings and accomplishments.

Concerning some of the comments above, all I can say is that it is not wise to judge a religion by its lowest members. If there was indeed an ISKCON brahmachari who was in a bar and begging for change, those are two violations of ISKCON principles right there. Therefore not a good representative of the movement. Likewise that sticker that was posted is not one that I have ever come across in relation to ISKCON before. There are many White Hindu organizations today, including other branches of Gaudiya Vaishnavism, that have nothing to do with ISKCON, so one shouldn't assume that any White person who claims to be a Hindu is a Hare Krishna.

The aggressive proselytization perpetrated by ISKCON is one of the factors that led to me leave the organization, but it is really no different in tone from that perpetrated by Christians and Muslims. Gaudiya Vaishnavism is a monotheistic philosophy and in some ways is very Abrahamic in its outlook, as in its desire to win converts. But for a Christian to condemn ISKCON for its proselytization practices (if indeed you are) would be extremely hypocritical.

Lastly, I'll say that ISKCON today is a large, international organization, and has many different branches that have evolved in different directions over the years. Some have deviated quite widely from the original philosophy. Some, such as the branch led by Bhakti Vikas Swami, are very conservative and have tried to stay as close to Prabhupada's and scriptural teachings as possible. So any statement one attempts to make about ISKCON as a whole will ultimately be incomplete.

\hfill

\texttt{John Morgan on 2013-06-19 at 08:46 said:}

``It's odd that an external enemy can be a spiritual friend, while some who claim to be close to us are in fact spiritual enemies. The latter are the worst kinds of enemy.''

I think this is a very profound and correct statement. However, I imagine that not merely the ability to appreciate a Shi'a philosopher is sufficient on its own to be considered a friend of traditionalism.

\hfill

\texttt{Jason-Adam on 2013-06-19 at 12:22 said:}

Mr Morgan :

I have never attended an ISKCON service or had any personal experience with the group at all, aside from reading a few of Prabhupada's books (which I rather enjoyed) so I am merely an observer of things here.

May I ask, are you sure that ISKCON offers a genuine initiation and is not akin to say a Blavatsky or a Vivekananda ? I am suspicious of ISKCON because among its adherents were rock stars, punk rockers, and other such gems of society... ..also there is this \url{http://prabhupada-krishna.co.uk/spiritual\_communism.htm} how can anyone be called traditional who uses words such as ``Vedic socialism''... ..also it seems as if Prabhupada ignores the idea of caste as something hereditary... ... ..

\hfill

\texttt{Cologero on 2013-06-19 at 12:46 said:}

The only ``genuine'' initiation in Hinduism, or better, Brahmanism, is the caste initiation. That is why outsiders cannot properly become Hindus. Yoga and Tantra groups do not discriminate on that basis, and probably on no basis.

\hfill

\texttt{Jason-Adam on 2013-06-19 at 15:02 said:}

So, Cologero, by definition that means any ``Hindu'' group that accepts Westerners is fraudulent ?

\hfill

\texttt{John Morgan on 2013-06-20 at 20:05 said:}

Cologero, in the Bhagavata Purana it states that anyone born in Kali-yuga is a shudra, regardless of whatever hereditary caste they belong to. (One can see this by the fact that the caste system in India today is in a complete shambles.) In the Gaudiya Vaishnava sampadaya, this is taken to mean that, during Kali-yuga, it is possible for those born in other castes, or even non-Indians, can become brahmanas, since brahmanas as exist in earlier yugas simply do not exist at this time.

\hfill

\texttt{John Morgan on 2013-06-20 at 20:34 said:}

I hasten to add, however, that in Gaudiya Vaishnavism, varanasrama (the caste hierarchy) is still considered important, but one acquires varna (caste) through behavior and learning rather than from birth, as in earlier yugas.

\hfill

\texttt{John Morgan on 2013-06-20 at 21:02 said:}

Jason-Adams, the only rock star I can think of who supported ISKCON in Prabhupada's time was George Harrison. Prabhupada always welcomed his support (among other things, Harrison purchased the Bhaktivedanta Manor in the UK for ISKCON, which is still one of its most important centers today), as he did from anyone who was genuinely interested, but Harrison never actually became a disciple of Prabhupada and he was never initiated. Prabhupada was an eminently practical man -- he welcomed support wherever he could get it, but never compromised his principles to do so. I imagine he considered it better for Harrison to become involved with Krishna on some level, even if he was unable to bring himself to embrace the austere life of a brahmachari or grihastha in that particular lifetime.

When Prabhupada spoke about ``spiritual Communism'' and ``Vedic socialism'' and so forth, what he meant was that he considered ISKCON a non-sectarian movement that encourages the increase of God-consciousness everywhere. I always considered this to be somewhat disingenuous, since ISKCON always maintains that Krishna is God and that any other form of God (Shiva, Ahura Mazda, Jesus, whatever) is only an incomplete version of Him. Still, it is true that Gaudiya Vaishnavas do not reject other religions, and encourage their growth for bringing people closer to God in some form, even if they regard GV as the most perfect form of religion. Prabhupada was certainly opposed to Communism and socialism in the usual sense, as he made clear in many of his writings and interviews. He was always a defender of varanasrama.

The idea of a valid initiation is that there is an unbroken line of initiations between gurus and disciples leading back to an initial divine revelation. In the Gaudiya Vaishnava sampradaya, the divine revelation was that of Lord Chaitanya, and there is most definitely an unbroken disciplic chain of succession from Chaitanya (and, in turn, Krsna) and Prabhupada. This can be found many places online. Blavatsky and Vivekananda could make no such claim.

\hfill

\texttt{Cologero on 2013-06-20 at 21:55 said:}

John Morgan, if you are trying to make a point beyond the confirmation of what we have written here many times over, I don't see it. By definition, ``shudra'' means one who does not have a spiritual initiation. In the process of degeneration of castes, the shudra wrest control of administration and rule, to the point that only the shudra and Brahmin exist. Ultimately, when all that remain are shudra, there is no longer any spiritual initiation. That, allegedly, is the situation in the West and, as you claim, may be in process even in India at this time. Nevertheless, at the current time, a Westerner cannot go to India and be initiated as a Brahmin, unless you know otherwise.

So it stills leaves the question, the answer to which we have been exploring, is how initiation is re-established when it exists nowhere, i.e., there are only shudras.

So the alternative, as you say, is through behavior and learning. That is what Gornahoor is encouraging our readers to do, sort of. (We don't think anyone should become a hare krishna).

As a corollary, any attempts at premature political solutions can only be the battle between one group of shudras (which are called ``common'' or ``vulgar'' in the West) and another, is simply futile.

\hfill

\texttt{John Morgan on 2013-06-21 at 05:17 said:}

Cologero, I'm not surprised that there are corollaries, given the reality of Tradition, although I am not giving you my personal opinions here, this is Gaudiya Vaishnava theology. So not original to me. As for brahmana initiations, thousands of people, including many Westerners, have been brahmana initiated in ISKCON, and in the Gaudiya Vaishnava sampradaya more broadly, and even more broadly in the other Vaishnava sampradayas, over the last half-century, and they are acknowledged as such by most Vedic authorities. Brahmana status, as exists in earlier yugas, is not possible at this time, even for Indians. So brahmana initiation now means something different than it does in other ages, but it is absolutely still going on and will be able to for at least another 9,500 years, as is confirmed in scripture. (Kali-yuga itself has about 430,000 years left to run, but eventually humans will degenerate to an animal level and become incapable of any spiritual impulse or activity.)

Personally I see Hare Krishna initiation as being just as valid for a Westerner as any other at this time, although it does create a wall of cultural separation between oneself and the rest of one's own people. But then again, if enlightenment and spiritual liberation is the primary goal, perhaps that shouldn't be important. Certainly Guenon would not have thought so. And many of the most seriously religious White individuals I have known were Hare Krishnas. All of the things that are being called for by this blog are certainly already being attempted, and in some cases achieved, by ISKCON.

\hfill

\texttt{Cologero on 2013-06-22 at 09:29 said:}

Curious, John Morgan. I presume, however, that no birth rite Hindu recognizes the caste change? Specifically, will the government recognize it, with their byzantine classification schemes? And an ISKCON initiate, moreover, cannot pass his initiation onto his son (unless the son has met the ISKCON requirements) ... is that true?

That entails becoming a Hindu, which is much more than the study of the six philosophical systems apart from any religious practice. Ironically, Venner on his blog rued Christianity and praised Hinduism for its ethnic closure. Obviously, he didn't know about ISKCON making Hinduism also universalistic.

\hfill

\texttt{John Morgan on 2013-06-22 at 10:07 said:}

Cologero, there are many hereditary Indian brahmanas in ISKCON today (some of them are among my friends), and of course they must acknowledge the hierarchy, which means that sannyasis of White ancestry or Indians born into lower varnas are their superiors, so no, it's not the case that all birth brahmanas reject this system. They have seen firsthand for themselves the wretched state that hereditary varnasrama is in today. I am certain that there are birth brahmanas who do not acknowledge ISKCON initiations, however. But I never encountered any hostility myself, even though I was an uninitiated neophyte. During my time with ISKCON in Mumbai, when I used to go out in my neighborhood dressed in dhoti-kurta and wearing tilak, I used to sometimes have Indians bow to me or who addressed me as ``pandit'' (or who shouted out ``Hare Krsna!'' when they saw me). So ``White Hindus'' have actually managed to garner respect even in India, in large part due to Prabhupada's efforts and the strict observances of his followers and due to the poor standards of most other Hindu organizations today (ISKCON is one of the most strongly conservative Hindu movements in India today). Anyway, even in earlier yugas, someone can lose his varna due to degenerate behavior, and it is sometimes possible for someone from a lower varna to acquire higher status. So it's never been entirely hereditary. The most important ways to determine varna, according to scripture, are through one's behavior and actions.

In terms of the Indian government today, there is nothing about caste in the Constitution at all, apart from the fact that varnasrama is abolished in favor of democracy, so no system of caste has political status in India today whatsoever, hereditary or otherwise. The BJP got in hot water recently because it insisted on having caste included as part of the last census.

And no, ISKCON initiations are not hereditary. One's sons must go through the same exact process as the father, if he wants to attain the same status.

Regarding Hinduism as an ethnically-based religion, it is said in the Vedic scriptures that anyone born outside of Bharat (India) is automatically less than a shudra (basically an animal), since only those born in India can attain full spiritual knowledge and escape samsara. It is also said that the scriptures cannot be understood in any language but Sanskrit. However, during Kali-yuga this is true of everyone, including Indians, so that is why it now becomes possible for non-Indians to become brahmanas. Until the 1930s, apart from a few rare converts, the phenomenon of White Hindus was quite unusual, which is surely where Guenon got this notion. There is a record of a Greek ambassador, Heliodorus, who converted to Vaishnavism in the fourth century BC. There was also some sannyasis, including a Gaudiya Vaishnava sannyasi, who travelled to the U.S. in the late 19th/early 20th century and made some converts there. But an earnest attempt to spread Gaudiya Vaishnavism outside India didn't begin until the 1930s, with the Gaudiya Math, and their efforts mostly met with failure until one of their students, A. C. Bhaktivedanta Prabhupada, went off to the U.S. in 1966, and then the rest of the world, and created ISKCON. Venner probably wasn't aware that the Hare Krishnas are a legitimate branch of the Vedic tree. But we shouldn't be too hard on him since, traditionally and until recently, that was indeed the case that Hinduism was an ethnically closed religion, and there are still branches of Hinduism that are closed to Westerners today.

\hfill

\texttt{Cologero on 2013-06-22 at 10:30 said:}

I thought that India had a system of scheduled castes and tribes, so I assumed caste had some legal status. I believe I read that an Indian court was trying to undo the stranglehold of the Brahmins on temple access. BTW, I wasn't being hard on Venner, since I had the same assumption. But the larger point (which the ISKCON philosophy implies) is that ethnicity cannot, in our time, be the primary factor; the fundamental problem in the West is the lack of a spiritual tradition, and population replacement is the symptom.

It raises an interesting point, assuming the ISKCON initiation is valid (and I presume it is, if it is done properly). In the west, there is a distinction between a valid and a licit ordination. A priestly ordination is valid if the administering bishop has apostolic succession. It is licit only if approved by the proper authorities. The analogous situation in the West is ``sedevacantism'', i.e., the position that those who are externally responsible for the priestly functions are not living up to their vocation. This is the same as ISKCON's claim that the outward Brahmins are not really Brahmins. So there are groups in the West who could conceivably perform the similar function. Unfortunately, they all fight over who has the purest doctrine and don't bother establishing communities of the type that you say ISKCON is forming in India.

\hfill

\texttt{John Morgan on 2013-06-22 at 11:12 said:}

Cologero, yes, there is such a thing as ``scheduled castes'' in India, by which is meant dalits and shudras. But this is only a form of ``affirmative action'' intended to undo the last vestiges of caste ``prejudice'' from India, not an attempt to integrate varnasrama into modern India. The modern Indian state is strongly Marxist in spirit. Even many erstwhile Hindu ``nationalists'' are opposed to caste. Caste may still plan some role in the Hindu civil courts -- the various religious communities in India each have their own courts for family and communal matters -- but I haven't studied the matter in detail.

I agree about the problem of spirituality as being greater than the ethnic one, even though I believe, along with Evola, in the importance of race as a traditional form.

And I also found ISKCON's notion of setting up farm communities as being an interesting one for members of any religious community to follow in the modern world.

I was satisfied that ISKCON's system of initiations is valid, although of course outsiders have challenged the legitimacy of ISKCON's system on various grounds over the years. And certainly there were many abuses, especially in the early years. But at least ideally it seemed like a good system to me, and especially ISKCON in India has really undergone a lot of positive changes over the last 20 years, and I imagine Indians will assume the forefront of the movement in coming years, as it seems largely dead in the water in the U.S. and Western Europe these days (not without some justification), apart from a few pockets here and there. As I wrote before, however, my problem with it was mainly that really embracing it fully essentially means giving up one's identity as a Westerner. But many see that as immaterial next to the spiritual goal.

As an aside, I'll mention that I've learned that in Goa there are Catholic brahmanas, Catholic kshatriyas, and so forth. This came about because, when the Portuguese first began converting the locals in the 16th century, the Church conducted negotiations with the Hindus, since the upper-caste Hindus didn't want to lose their status in order to become Catholic. And apparently there are still some specifically Goan rituals in the Church that can only be performed by brahmanas and kshatriyas today.

\hfill

\end{sffamily}
\end{footnotesize}

